% Artigo no template SBC. Requer sbc-template.sty e sbc.bst (baixar do site da SBC:
% https://www.sbc.org.br/documentos-da-sbc/category/169-templates-para-artigos-e-capitulos-de-livros).
% Compilar: pdflatex artigo; bibtex artigo; pdflatex artigo; pdflatex artigo
% Marcações TODO(...) indicam trechos que dependem do trabalho de outro integrante.

\documentclass[12pt]{article}

\usepackage{sbc-template}
\usepackage{graphicx,url}
\usepackage[brazil]{babel}
\usepackage[utf8]{inputenc}
\usepackage{booktabs}

\sloppy

\title{GraphRAG e KAG sobre um Catálogo de Filmes e Séries: \\
Recuperação em Grafo de Propriedades para Reduzir Alucinações em LLMs}

\author{Ewellyn Maria\inst{1}, Gustavo Palmeira\inst{1}, Rafael Cirne\inst{1},\\
Julliane Di Paula\inst{1}, Felipe Souto Maior\inst{1}, Arthur Xavier\inst{1}}

\address{UNIPÊ -- Centro Universitário de João Pessoa\\
BR-230 -- Água Fria, João Pessoa -- PB, 58053-000
\email{\{ewellyn.andrade, gustavo.martinez, rafael.medeiros,}\\
{\footnotesize\texttt{julliane.xavier, felipe.mendes, acavalcante\}@cs.unipe.edu.br}}}

\begin{document}

\maketitle

\begin{abstract}
Large language models answer fluently even when they lack the facts, which produces
hallucinations. This paper presents a GraphRAG/KAG prototype that stores a catalog of movies and TV
series in a Neo4j property graph and answers multi-hop questions from retrieved, traceable evidence.
We compare three variants---an LLM without retrieval, vector RAG and hybrid GraphRAG---on
17 questions whose gold answers are derived from the catalog itself, including questions that
the catalog cannot answer. [TODO: resumir resultado numérico após a execução final.]
\end{abstract}

\begin{resumo}
Modelos de linguagem respondem com fluência mesmo sem conhecer os fatos, o que gera alucinações.
Este trabalho apresenta um protótipo de GraphRAG/KAG que armazena um catálogo de filmes e séries
em um grafo de propriedades no Neo4j e responde perguntas multi-hop a partir de evidências
recuperadas e rastreáveis. Comparamos três variantes (LLM sem recuperação, RAG vetorial e GraphRAG
híbrido) em 17 perguntas cujo gabarito é derivado do próprio catálogo, incluindo perguntas que o
catálogo não consegue responder. [TODO: resumir o resultado numérico após a execução final.]
\end{resumo}

\section{Introdução}

Modelos de linguagem de grande porte (LLMs) geram texto plausível independentemente de possuírem a
informação necessária, e esse comportamento é conhecido como alucinação~\cite{ji2023survey}. A
geração aumentada por recuperação (RAG) reduz o problema ao fornecer ao modelo trechos de texto
relevantes~\cite{lewis2020rag}, mas a busca por similaridade vetorial recupera trechos isolados e tem
dificuldade com perguntas que exigem combinar fatos espalhados em vários documentos, como
\emph{``quais atores atuaram em séries criadas por quem criou Better Call Saul?''}.

Abordagens baseadas em grafos de conhecimento atacam essa limitação ao representar entidades e
relações explicitamente, permitindo percorrer múltiplos saltos~\cite{edge2024graphrag,liang2024kag}.
Este trabalho constrói um protótipo acadêmico de ponta a ponta e responde à pergunta de pesquisa:
\textbf{a recuperação em grafo reduz a alucinação de um LLM em relação ao mesmo LLM sem recuperação e
ao RAG puramente vetorial?}

As contribuições são: (i) um pipeline Text-to-Graph sobre um catálogo de filmes e séries com
persistência em Neo4j; (ii) recuperação híbrida (vetor, texto completo e caminho no grafo) com
resposta fundamentada e Cypher exibido; e (iii) um protocolo de avaliação cujo gabarito é derivado
do catálogo e que inclui perguntas sem resposta, para medir se o sistema se abstém em vez de inventar.

\section{Fundamentação Teórica}

\subsection{Alucinação e RAG}
Alucinação é a geração de conteúdo não sustentado pela fonte ou pelo mundo~\cite{ji2023survey}. Em
sistemas RAG~\cite{lewis2020rag}, a fonte é o conjunto de trechos recuperados, e a fidelidade a ela
(\emph{grounding}) é a propriedade que buscamos medir.

\subsection{Grafos de propriedades e Cypher}
Em um grafo de propriedades rotulado (LPG), nós e relações possuem rótulos e propriedades. O Neo4j
armazena esse modelo e é consultado com Cypher~\cite{francis2018cypher}. Perguntas multi-hop
correspondem a padrões de caminho, por exemplo
$(Person)\!-\![:ACTED\_IN]\!\rightarrow\!(Movie)\!\leftarrow\![:DIRECTED]\!-\!(Person)$.

\subsection{GraphRAG e KAG}
O GraphRAG~\cite{edge2024graphrag} constrói um grafo a partir de texto e usa detecção de
comunidades e resumos para responder perguntas globais; o KAG~\cite{liang2024kag} integra grafo de
conhecimento e geração com raciocínio guiado por lógica. Algoritmos como PageRank~\cite{page1999pagerank}
e Leiden~\cite{traag2019leiden} podem enriquecer o contexto com a importância e a comunidade dos nós.
% TODO(Pessoa 3): confirmar quais algoritmos entraram na versão final.

\section{Arquitetura do Sistema}

O fluxo implementado é: fonte $\rightarrow$ documento $\rightarrow$ chunk $\rightarrow$
entidades e relações $\rightarrow$ Neo4j $\rightarrow$ recuperação híbrida $\rightarrow$ contexto
fundamentado $\rightarrow$ interface.

\begin{itemize}
  \item \textbf{Dados:} catálogo de 9 títulos (5 filmes e 4 séries) com diretores/criadores, elenco,
        gêneros e estúdio, conferidos na Wikipedia (CC BY-SA 4.0).
  \item \textbf{Text-to-Graph:} nós \texttt{Movie}, \texttt{Series}, \texttt{Person}, \texttt{Genre},
        \texttt{Studio}, \texttt{Document}, \texttt{Chunk} e \texttt{Entity}; relações \texttt{DIRECTED},
        \texttt{CREATED}, \texttt{ACTED\_IN}, \texttt{HAS\_GENRE}, \texttt{PRODUCED\_BY},
        \texttt{HAS\_CHUNK}, \texttt{MENTIONS} e \texttt{RELATED\_TO}.
        % TODO(Pessoa 2): NER/RE final e ligação Document--Movie/Series.
  \item \textbf{Recuperação:} índice vetorial (cosseno, 384 dimensões), índice de texto completo e
        pontuação por caminho mais curto, combinados em um escore híbrido.
        % TODO(Pessoas 1 e 3): embeddings semânticos, geração de Cypher por LLM e multi-hop real.
  \item \textbf{Geração:} LLM local (Ollama) com as evidências como único contexto; sem Ollama,
        resposta por \emph{template}.
\end{itemize}

\section{Metodologia de Avaliação}

\subsection{Conjunto de perguntas}
São 17 perguntas em cinco categorias: 1 salto (3), 2 saltos (6), 3 saltos (2), agregação (1) e
\emph{sem resposta no catálogo} (4). O gabarito de cada pergunta (entidades esperadas e documentos de
suporte) é \textbf{calculado a partir de \texttt{titles.json}}, e não escrito à mão, o que evita
inconsistência e permite regerar o conjunto se o catálogo crescer. As perguntas sem resposta
(por exemplo, um filme que não está no catálogo) verificam se o sistema admite não saber; a geração
do conjunto falha se o catálogo passar a conter a resposta, para que essas perguntas continuem válidas.

\subsection{Variantes comparadas}
\begin{enumerate}
  \item \textbf{LLM sem recuperação:} o mesmo modelo recebe apenas a pergunta, com prompt neutro.
  \item \textbf{RAG vetorial:} os 5 documentos mais próximos pelo índice vetorial viram contexto.
  \item \textbf{GraphRAG híbrido:} contexto da recuperação híbrida com caminho no grafo.
  \item \textbf{LLM externo sem recuperação} (opcional): respostas fornecidas por arquivo.
\end{enumerate}
O modelo gerador e o prompt de fundamentação são idênticos nas variantes 2 e 3; apenas o contexto muda.

\subsection{Métricas}
\begin{itemize}
  \item \textbf{Recall de entidades:} fração das entidades do gabarito citadas na resposta.
  \item \textbf{Acerto:} perguntas respondíveis com recall total; perguntas sem resposta com abstenção
        e sem entidades inventadas.
  \item \textbf{Alucinação:} resposta que cita nome próprio ausente do catálogo e da pergunta, ou que
        responde com segurança uma pergunta sem resposta no catálogo.
  \item \textbf{Avaliação humana:} correção, relevância, completude e fundamentação (escala 1--5)
        em planilha-modelo, para validar a métrica automática.
\end{itemize}
A alucinação é medida em relação ao \emph{catálogo}, não ao mundo real: uma resposta verdadeira fora
do catálogo conta como não fundamentada. A detecção de entidades é heurística (nomes próprios com
inicial maiúscula) e subestima casos em que uma palavra única inicia a frase.

\section{Resultados}

% TODO: rodar `python -m scripts.evaluate_hallucination` com o pipeline final e colar a tabela.
\begin{table}[ht]
\centering
\caption{Resultados gerais (preencher após a execução final).}
\label{tab:geral}
\begin{tabular}{lrrrr}
\toprule
Variante & Acerto & Recall médio & Alucinação & Abstenção correta \\
\midrule
LLM sem recuperação & -- & -- & -- & -- \\
RAG vetorial        & -- & -- & -- & -- \\
GraphRAG híbrido    & -- & -- & -- & -- \\
\bottomrule
\end{tabular}
\end{table}

% TODO: tabela por categoria (1, 2 e 3 saltos, agregação, sem resposta) e análise de erros.

\section{Discussão e Limitações}

Os resultados devem ser lidos com cautela. O catálogo tem apenas 9 títulos, de modo que o
experimento demonstra o mecanismo, mas não permite conclusões estatísticas. Os embeddings usados são
por \emph{hashing} determinístico e não capturam semântica, o que favorece a busca por termos exatos.
O gerador é um modelo pequeno executado localmente, e a métrica automática depende de heurística de
extração de entidades. Por fim, o baseline sem recuperação tem vantagem injusta de conhecimento
paramétrico sobre os filmes reais, mas é penalizado por usar fatos fora do catálogo.
% TODO: revisar esta seção com os números finais e as limitações vigentes dos demais módulos.

\section{Conclusão}

% TODO: conclusão após os resultados. Resumir (i) o que o grafo acrescentou em perguntas multi-hop,
% (ii) o comportamento de abstenção, e (iii) trabalhos futuros (catálogo maior, embeddings
% semânticos, citações e comunidades, avaliação humana ampliada).

\bibliographystyle{sbc}
\bibliography{Referencias/referencias}

\end{document}
